\documentclass[conference]{IEEEtran}

\usepackage{cite}
\usepackage{amsmath,amssymb}
\usepackage{booktabs}
\usepackage{array}
\usepackage{url}
\usepackage{hyperref}

\title{A Preliminary Study on Scientific Claim Verification to Identify Citation Drift in AI-Generated Literature Summaries}

\author{
\IEEEauthorblockN{[Author Name]}
\IEEEauthorblockA{
[Department / School]\\
[Institution]\\
[City, Country]\\
[Email Address]
}
}

\begin{document}
\maketitle

\begin{abstract}
We have previously identified that literature summaries generated by some Artificial Intelligence (AI) systems may include incomplete or incorrect scientific claims connected to misleading or incorrect citations. In this preliminary study, we develop and evaluate a basic method to identify ``citation drift'' in literature summaries created by AI systems. To do this, we use a three-way natural language inference (NLI) cross-encoder model, \texttt{cross-encoder/nli-deberta-v3-base}, to predict contradiction, entailment, or neutral relationships between claim--evidence pairs using the development set of the SciFact dataset ($n=338$). Predictions of entailment are mapped to SUPPORT, while predictions of contradiction are mapped to CONTRADICT; all other predictions are classified as NEUTRAL. When NEUTRAL predictions are considered incorrect, the overall accuracy of the model is 45.56\%. The model makes a non-neutral prediction (either SUPPORT or CONTRADICT) for 50.30\% of the examples, and 90.59\% of these predictions are correct (macro F1-score = 0.91). In summary, when the model predicts either SUPPORT or CONTRADICT, it is accurate on this subset; however, it does not make a directional prediction for nearly half of the examples. This is a preliminary component-level baseline and does not establish citation-drift detection performance.
\end{abstract}

\begin{IEEEkeywords}
Citation drift, scientific claim verification, natural language inference, literature summarization, SciFact, factual consistency.
\end{IEEEkeywords}

\section{Introduction}
The use of large language models (LLMs) to summarize scientific papers is becoming more prevalent. Along with summaries of the methods and results of the studies included in the papers, users of LLMs may also expect summaries to describe potential limitations of the studies, as well as to provide a synthesis of related studies conducted in the same research area. However, when synthesizing information across several studies, it is not uncommon for the summary to contain inaccuracies. For example, a summary may include or exclude qualifiers, state a result differently, or describe it more vaguely than in the original study. Such inconsistencies provide the motivation for this work to study citation drift.

Reliably determining the accuracy of claims in scientific literature is closely related to natural language inference (NLI). In this setting, the summarized claim is treated as the hypothesis and the supporting evidence as the premise. However, compared with other fields, the use of scientific terminology and evidence to support claims incorporated in the text greatly increases the complexity of the task.

We present a preliminary study using an NLI model to classify scientific claim--evidence pairs from the SciFact dataset.

The experiment has three main aspects. First, we use a pretrained DeBERTa-v3 NLI cross-encoder to assess pairs of scientific claims and evidence. Second, we take into consideration the three classes associated with the model output, in contrast to the two gold (human-annotated) labels used in this evaluation task. Third, we report both the overall accuracy of predictions and conditional accuracy, allowing the reader to assess the contribution of NEUTRAL predictions.

\section{Related Work}

\subsection{Scientific Claim Verification}
Scientific claim verification gives researchers the means to automatically check whether evidence in the research literature supports or refutes a given scientific claim. Wadden \emph{et al.} created the SciFact dataset, which contains scientific claims paired with evidence from research paper abstracts, together with support/refutation labels and rationales \cite{wadden2020fact}. We use this dataset for claim--evidence pair evaluation. In particular, we do not consider SciFact to contain summaries generated by natural language generation (NLG) systems.

\subsection{Natural Language Inference}
In natural language inference, reasoning between pairs of sentences is formulated as a three-way classification problem: deciding whether the relationship is entailment, contradiction, or neutral. The SNLI and MultiNLI datasets are examples of resources developed for this form of inference \cite{bowman2015large,williams2018broad}.

\subsection{Previous Language Models}
One method for language representation learning uses disentangled attention along with an enhanced mask decoder, as proposed in DeBERTa. A newer version of this model family, DeBERTa-v3, uses replaced-token detection along with gradient-disentangled embedding sharing. We employ the pretrained model \texttt{cross-encoder/nli-deberta-v3-base} for our NLI task; no further training is performed \cite{he2021deberta,he2023debertav3}.

\subsection{Positioning of This Study}
We are in the early stages of applying this methodology and therefore present a simple baseline solution for evaluating scientific claims associated with automatically generated summaries.

Our current implementation does not include automated citation links based on generated summaries or manual annotations to evaluate citation drift. Such an evaluation would require human assessment of each individual atomic claim together with the text supporting that claim.

\section{Method}

\subsection{Formulation of the Task}
In each example, we are given a scientific claim $c$ and an evidence text $e$. We provide the evidence text $e$ to the NLI model as the premise and the scientific claim $c$ as the hypothesis. The model then generates one of three classifications: CONTRADICTION, ENTAILMENT, or NEUTRAL.

For comparison with the labels provided in the SciFact dataset, we treat ENTAILMENT as SUPPORT and CONTRADICTION as CONTRADICT.

\subsection{NLI Inference}
We use the cross-encoder model \texttt{cross-encoder/nli-deberta-v3-base} to score pairs of evidence and claims. The model output indices are mapped as follows: 0 = contradiction, 1 = entailment, and 2 = neutral. Based on this output, we predict the class for each evidence--claim pair. We do not apply any further training or optimization to this model using the examples available in the SciFact development set.

\subsection{Evaluation Measures}
We calculate overall accuracy over all 338 examples in the development set. We also calculate the coverage of the set of examples for which the model returns either SUPPORT or CONTRADICT. Finally, we compute conditional precision, recall, and F1-score over this subset of examples for which the model returns either SUPPORT or CONTRADICT.

\section{Experimental Setup}

\subsection{Dataset}
We use SciFact to obtain claims, evidence rationales, and abstracts. We create 338 claim--evidence pairs using the development split of SciFact. Of these, 216 examples are SUPPORT and 122 are CONTRADICT according to the gold labels.

\subsection{Reproducibility Artifacts}
The experimental artifacts include a CSV containing claim--evidence pairs and predictions, a Colab notebook, and the printed evaluation outputs. For reproducibility, the project should retain the notebook, dataset release information, model identifier or revision, and package versions.

\section{Results and Discussion}

\subsection{Prediction Distribution}
The model produces 76 SUPPORT predictions, 94 CONTRADICT predictions, and 168 NEUTRAL predictions. Therefore, it makes a non-neutral prediction for 170 out of 338 examples, giving coverage of 50.30\%.

\begin{table}[t]
\centering
\caption{Prediction distribution over 338 claim--evidence pairs.}
\label{tab:prediction-distribution}
\begin{tabular}{lr}
\toprule
\textbf{Predicted class} & \textbf{Count}\\
\midrule
SUPPORT & 76\\
CONTRADICT & 94\\
NEUTRAL & 168\\
\midrule
Total & 338\\
\bottomrule
\end{tabular}
\end{table}

\subsection{Overall Evaluation}
The full confusion matrix, with gold labels as rows and predicted labels as columns, is shown in Table~\ref{tab:full-confusion}.

\begin{table}[t]
\centering
\caption{Full confusion matrix (gold labels in rows and predictions in columns).}
\label{tab:full-confusion}
\begin{tabular}{lrrr}
\toprule
\textbf{Gold label} & \textbf{SUPPORT} & \textbf{CONTRADICT} & \textbf{NEUTRAL}\\
\midrule
SUPPORT & 73 & 13 & 130\\
CONTRADICT & 3 & 81 & 38\\
\bottomrule
\end{tabular}
\end{table}

The overall accuracy is 45.56\%, calculated as 154 correct predictions out of 338. NEUTRAL predictions are counted as incorrect because the gold-label set contains only SUPPORT and CONTRADICT.

\subsection{Conditional Evaluation}
When NEUTRAL predictions are excluded, 170 examples remain. The model correctly classifies 154 of them, giving conditional accuracy of 90.59\%. The conditional confusion matrix is shown in Table~\ref{tab:conditional-confusion}.

\begin{table}[t]
\centering
\caption{Conditional confusion matrix for non-neutral predictions.}
\label{tab:conditional-confusion}
\begin{tabular}{lrr}
\toprule
\textbf{Gold label} & \textbf{SUPPORT} & \textbf{CONTRADICT}\\
\midrule
SUPPORT & 73 & 13\\
CONTRADICT & 3 & 81\\
\bottomrule
\end{tabular}
\end{table}

For this subset, the reported classification metrics are shown in Table~\ref{tab:conditional-metrics}.

\begin{table}[t]
\centering
\caption{Classification metrics on the 170 non-neutral predictions.}
\label{tab:conditional-metrics}
\begin{tabular}{lrrrr}
\toprule
\textbf{Class} & \textbf{Precision} & \textbf{Recall} & \textbf{F1} & \textbf{Support}\\
\midrule
CONTRADICT & 0.86 & 0.96 & 0.91 & 84\\
SUPPORT & 0.96 & 0.85 & 0.90 & 86\\
Macro average & 0.91 & 0.91 & 0.91 & 170\\
Weighted average & 0.92 & 0.91 & 0.91 & 170\\
\bottomrule
\end{tabular}
\end{table}

\subsection{Discussion}
The main observation is a trade-off between prediction coverage and conditional accuracy. The model makes a non-neutral prediction for approximately half of the examples. Within this subset, its predictions agree with the gold labels in 90.59\% of cases. However, the conditional score applies only to the 170 non-neutral predictions. Reporting it without the 50.30\% coverage would overstate performance across the full dataset.

The 168 NEUTRAL predictions may reflect uncertainty or a mismatch between the model's NLI decision boundary and the way evidence is presented. The current measurements do not establish the cause. The results suggest that a pretrained NLI model can provide a useful starting point for scientific claim verification, but further analysis is required before using it as a citation-verification component.

\section{Limitations}
This study has several limitations. First, SciFact is a scientific claim-verification dataset, not a dataset of AI-generated literature summaries and their citations. Second, the experiment uses the development split and does not report an independent test-set evaluation. Third, the gold annotations contain two classes, while the model produces three classes. NEUTRAL predictions are counted as incorrect for overall accuracy, and this treatment affects the interpretation of the result. Fourth, conditional accuracy applies only to non-neutral predictions and must be reported with coverage. Fifth, the exact runtime configuration, package versions, model revision, and some dataset-release details are not included in the available experiment record. Finally, this experiment does not evaluate retrieval quality, evidence aggregation across multiple sources, atomic claim decomposition, or human annotator agreement.

\section{Conclusion and Future Work}
This paper presents a preliminary scientific claim-verification baseline using a pretrained DeBERTa-v3 NLI cross-encoder on 338 SciFact development claim--evidence pairs. The model achieves 45.56\% overall accuracy when NEUTRAL predictions are counted as incorrect. It produces non-neutral predictions for 50.30\% of the examples and achieves 90.59\% conditional accuracy on that subset. These results show that coverage should be reported alongside conditional performance when a model has a NEUTRAL output that is not represented in the gold-label set.

The experiment is a component-level baseline relevant to citation-support assessment. It is not a direct measurement of citation drift in AI-generated literature summaries. Future work should construct a benchmark from AI-generated summaries, break summaries into atomic claims, connect each claim to its cited source, retrieve relevant source passages, and obtain human-verified support labels. The evaluation should also distinguish genuinely unsupported claims from cases where relevant evidence is unavailable or not retrieved.

\bibliographystyle{IEEEtran}
\bibliography{references}

\end{document}
